%!TEX program = uplatex
\RequirePackage{plautopatch}
\documentclass[uplatex,dvipdfmx]{jlreq}
\usepackage{amsmath,amssymb}

\pagestyle{empty}

\begin{document}

%======================================================================
% 表紙
%======================================================================
\begin{center}
  \vspace*{20pt}
  {\Huge\bfseries 幾何学者は物理学から何を学んだか}

  \vspace{1pc}
  {\Large 概要}
\end{center}

\vfill

\newpage

%======================================================================
% 講演１：深谷賢治「数学と物理のこの20年」
%======================================================================
\noindent
{\large\bfseries 数学と物理のこの20年}

\bigskip
\hspace*{\fill}{\large\bfseries 深谷　賢治（京大・理）}

\bigskip

幾何学の世界に理論物理学の影響が顕著になってかなりの年月がたちました。
これは、同時に理論物理学で、大域的な幾何学の見地の役割の重要性が
次第に明らかになってきた歴史でもあります。

例えば、アノーマリ、超対称性、ファインマン経路積分、超弦理論、
位相的場の理論、共形場の理論などの、理論物理学に属するトピックと、
Atiyah--Singer の指数定理、さまざまな非線形方程式の解のモジュライ空間、
3次元多様体や4次元多様体の不変量、Hyper Kähler Geometry、
大域シンプレクティック幾何学（Symplectic topology）などの多くの話題が、
相互に深くかかわりながら発達しました。
その流れを振り返りながら、いろいろな考えをみていきたいと思います。

これらのうちで、物理学に属するものの多くは、
それが発見された頃に、物理学者に話を聞かせてもらった事柄です。
しかし、私なども、当時、未消化なまますませてしまったことが多かったように思います。
また、私などの理解では、それぞれ、断片的に聞こえてくるだけで、
多くの物理学と幾何学の関わりが、全体としてどんなストーリーになっているか、
なかなか見えませんでした。
いまでも、私が、きちっと全体像を把握しているわけではとうていありませんが、
この機会に、少しでもまとまった形でお話しできたらと思います。

理論物理学に属する事柄に関して、私は、当然、専門家より理解・知識が不足しているわけですが、
他面、数学を専門とする人間にとって、なにが分かりずらいか、
また、どんな言葉遣いが自然であるか、にはより通じていると思います。
数学者が物理学者の言うことを、いつも聞くだけであると、結局未消化に終わる、
というのが、私自身の経験から来る、過去の教訓であるように思えますので、
数学者の言葉で語る、という試みを、出来るだけしてみようと思います。

お話する題材をなににするかは、まだ決めかねていることも多いのですが、
そうできたらいいと考えているのは、例えば、次のような点です。
\begin{itemize}
  \item 無限次元にかかわる考察から、有限の量が引き出せる。
        そして、それを調べることが出来るのは、どのようにしてか。
  \item 理論物理学で、大域的な幾何学、例えば空間の位相幾何学的性質が、
        どうして、また、どのように現れるか。
  \item 理論物理学のどのような概念が、数学のどのような概念とかかわっているのか。
\end{itemize}

勿論、このようなことを一般的に話す、のが予定ではなく、
それを示している具体的な例を話してみたい、と考えています。

\newpage

%======================================================================
% 講演２：古田幹雄「Donaldson 不変量の構造定理への4つのアプローチ」
%======================================================================
\noindent
{\large\bfseries Donaldson 不変量の構造定理への4つのアプローチ}

\bigskip
\hspace*{\fill}{\large\bfseries 古田　幹雄（京大・数理研）}

\bigskip

数学と物理とが異なる角度からひとつの現象を照らしだす、そうした一例の紹介が
今回の話の目標です。

4次元多様体の Donaldson 不変量は、（偶数次元ホモロジー群上の）同次多項式の族です。
すなわち、4次元多様体上に $SU(2)$ 束が与えられると、
その上の $SU(2)$ 接続全体の空間を出発点としてある手続きを経ることにより、
ひとつの同次多項式を定義できます。
$SU(2)$ 束を動かすと、無限個の同次多項式が得られますが、それらは互いに独立ではなく、
ある有限個の2次コホモロジー要素を用いて、一挙に表示できる事実が知られています。
これらの有限個の要素は4次元多様体の「basic class」とよばれ、
複素多様体における「canonical class」に比すことができます。

この事実の「証明／説明」には、現在ほぼ4通りのアプローチがあります。
2つは数学から、2つは物理からです。

数学的アプローチで問題となるのは「相異なるトポロジーを持つ $SU(2)$ 束を同時に扱うこと」です。
物理において、これはあまり問題になりません。
「トポロジーをも動かしてすべての可能な $SU(2)$ 接続に対して和をとる」のが出発点だからです。
物理で問題になるのは「真空のモジュライの構造は何か」です。

物理の議論は、過去の物理における経験の蓄積に基づいており、
講演者はそれを理解しているとはいいがたいのですが、
可能な範囲で説明を試みる予定です。

4つのアプローチとは以下の通りです。

\begin{description}
  \item[(1) 数学から（Kronheimer--Mrowka）]\mbox{}\\
        4次元多様体の内に曲面をひとつ固定し、その曲面に沿って特異性をもつ $SU(2)$ 束を導入する。
        それを利用して相異なるトポロジーを持つ2つの $SU(2)$ 束を1パラメータ族で結ぶ。

  \item[(2) 数学から（Pidstrigach--Tyurin）]\mbox{}\\
        $SU(2)$ 接続だけでなく、ベクトル束に値を持つスピノルをも考える。
        それを利用して「$SU(2)$ 接続の考察」を
        「$U(1)$ 接続と直線束に値をもつスピノルの考察」に帰着する。
        相異なる $SU(2)$ 束は「バブル」を通じて現れる。

  \item[(3) 物理から（Witten）]\mbox{}\\
        4次元多様体が Kähler 曲面であり、0でない正則2形式を持つと仮定する。
        それを利用して理論を変形し、「質量ギャップ」を持たせると、真空は有限個になる。
        真空のモジュライの決定は正則2形式の零点集合の近傍の考察に帰着する。

  \item[(4) 物理から（Seiberg--Witten）]\mbox{}\\
        $SU(2)$ 束の構造群を $U(1)$ まで簡約して $U(1)$ 接続の考察に帰着しようと考える。
        この簡約の「障害」は真空のモジュライ上では特異点として現れる。
        Donaldson 不変量への寄与は、この「障害」すなわち「モノポール」のみからやってくると主張される。
        この「障害」を「双対」で考察すると、
        「$U(1)$ 接続と直線束に値をもつスピノルの考察」に帰着する。
\end{description}

\end{document}
